
\documentclass{article}
\usepackage{colortbl}
\usepackage{makecell}
\usepackage{multirow}
\usepackage{supertabular}

\begin{document}

\newcounter{utterance}

\centering \large Interaction Transcript for game `cladder', experiment `full\_v1.5\_default', episode 4908 with Qwen3.6{-}27B.
\vspace{24pt}

{ \footnotesize  \setcounter{utterance}{1}
\setlength{\tabcolsep}{0pt}
\begin{supertabular}{c@{$\;$}|p{.15\linewidth}@{}p{.15\linewidth}p{.15\linewidth}p{.15\linewidth}p{.15\linewidth}p{.15\linewidth}}
    \# & \multicolumn{2}{c}{Player} && \multicolumn{2}{c}{Game Master} \\
    \hline

    \theutterance \stepcounter{utterance}  
    & & & \multicolumn{4}{p{0.6\linewidth}}{
        \cellcolor[rgb]{0.9,0.9,0.9}{
            \makecell[{{p{\linewidth}}}]{
                \texttt{\tiny{[P1$\langle$GM]}}
                \texttt{You are an expert in causal inference. The following question is not a typical commonsense query, but rather a meticulously designed question created by a professor specializing in causal inference, intended to assess the students' mastery of the course content.} \\
\texttt{Imagine a self{-}contained, hypothetical world with only the following conditions, and without any unmentioned factors or causal relationships: My decision has a direct effect on penguin mood and curly hair. Penguin mood has a direct effect on curly hair. For those who choose to take the stairs, the probability of curly hair is 70\%. For those who choose to take the elevator, the probability of curly hair is 17\%. For those who choose to take the elevator, would it be less likely to see curly hair if I had taken the stairs?} \\
\texttt{Start your answer with 'Yes' or 'No', followed by additional reasoning or evidence to support your explanation.} \\
            }
        }
    }
     \\ \\

    \theutterance \stepcounter{utterance}  
    & \multicolumn{4}{p{0.6\linewidth}}{
        \cellcolor[rgb]{0.561,0.737,0.561}{
            \makecell[{{p{\linewidth}}}]{
                \texttt{\tiny{[P1$\rangle$GM]}}
                \texttt{Yes} \\
\\ 
\texttt{**Reasoning:**} \\
\\ 
\texttt{To answer this question, we must distinguish between the **observed probability** (marginal probability) and the **counterfactual probability** (individual causal effect).} \\
\\ 
\texttt{1.  **Analyze the Causal Structure:**} \\
\texttt{    The problem states:} \\
\texttt{    *   Decision (\$D\$) \$\textbackslash{}rightarrow\$ Penguin Mood (\$M\$)} \\
\texttt{    *   Decision (\$D\$) \$\textbackslash{}rightarrow\$ Curly Hair (\$H\$)} \\
\texttt{    *   Penguin Mood (\$M\$) \$\textbackslash{}rightarrow\$ Curly Hair (\$H\$)} \\
\\ 
\texttt{    This implies that \$D\$ affects \$H\$ through two pathways: a direct effect and an indirect effect mediated by \$M\$.} \\
\\ 
\texttt{2.  **Analyze the Observed Data:**} \\
\texttt{    *   \$P(H=1 | D=\textbackslash{}text\{stairs\}) = 0.70\$} \\
\texttt{    *   \$P(H=1 | D=\textbackslash{}text\{elevator\}) = 0.17\$} \\
\\ 
\texttt{    These are average probabilities across the population. They tell us that, on average, people who take the stairs have curly hair 70\% of the time, and people who take the elevator have curly hair 17\% of the time.} \\
\\ 
\texttt{3.  **Analyze the Counterfactual Question:**} \\
\texttt{    The question asks: *"For those who choose to take the elevator, would it be less likely to see curly hair if I had taken the stairs?"*} \\
            }
        }
    }
    & & \\ \\

    \theutterance \stepcounter{utterance}  
    & & & \multicolumn{2}{p{0.3\linewidth}}{
        \cellcolor[rgb]{0.9,0.9,0.9}{
            \makecell[{{p{\linewidth}}}]{
                \texttt{\tiny{[GM$|$GM]}}
                \texttt{Yes} \\
            }
        }
    }
    & & \\ \\

    \theutterance \stepcounter{utterance}  
    & & & \multicolumn{2}{p{0.3\linewidth}}{
        \cellcolor[rgb]{0.9,0.9,0.9}{
            \makecell[{{p{\linewidth}}}]{
                \texttt{\tiny{[GM$|$GM]}}
                \texttt{no} \\
            }
        }
    }
    & & \\ \\

    \theutterance \stepcounter{utterance}  
    & & & \multicolumn{2}{p{0.3\linewidth}}{
        \cellcolor[rgb]{0.9,0.9,0.9}{
            \makecell[{{p{\linewidth}}}]{
                \texttt{\tiny{[GM$|$GM]}}
                \texttt{game\_result = LOSE} \\
            }
        }
    }
    & & \\ \\

\end{supertabular}
}

\end{document}
